% ============================================================
%  USER STORIES — Projet PFE
%  Langue : Français
%  NE PAS COMPILER — fichier source LaTeX uniquement
% ============================================================
\documentclass[12pt,a4paper]{article}

% ── Encodage & langue ──────────────────────────────────────
\usepackage[utf8]{inputenc}
\usepackage[T1]{fontenc}
\usepackage[french]{babel}

% ── Mise en page ───────────────────────────────────────────
\usepackage[
  top=2cm, bottom=2cm,
  left=2cm, right=2cm
]{geometry}
\usepackage{parskip}            % Espacement entre paragraphes
\usepackage{setspace}           % Interligne
\onehalfspacing

% ── Polices ────────────────────────────────────────────────
\usepackage{lmodern}
\usepackage{amssymb} % Pour \blacktriangleright

% ── Couleurs ───────────────────────────────────────────────
\usepackage{array} % DOIT ETRE AVANT xcolor[table] pour colortbl
\usepackage[table,dvipsnames,svgnames]{xcolor}

% Palette personnalisée
\definecolor{chefBlue}{HTML}{1B4F72}
\definecolor{chefBg}{HTML}{D6EAF8}
\definecolor{teacherGreen}{HTML}{1E8449}
\definecolor{teacherBg}{HTML}{D5F5E3}
\definecolor{studentOrange}{HTML}{D35400}
\definecolor{studentBg}{HTML}{FDEBD0}
\definecolor{headerBg}{HTML}{2C3E50}
\definecolor{accentGold}{HTML}{F39C12}
\definecolor{lightGray}{HTML}{F2F3F4}
\definecolor{darkText}{HTML}{2C3E50}
\definecolor{routeGray}{HTML}{7F8C8D}

% ── Tableaux ───────────────────────────────────────────────
\usepackage{tabularx}
\usepackage{booktabs}
\usepackage{multirow}
\usepackage{makecell}
\renewcommand{\arraystretch}{1.4}

% ── Listes ─────────────────────────────────────────────────
\usepackage{enumitem}

% ── Symboles (pifont) ──────────────────────────────────────
\usepackage{pifont}

% ── Comparaison de chaînes (etoolbox) ──────────────────────
\usepackage{etoolbox}

% ── Boîtes colorées (tcolorbox) ────────────────────────────
\usepackage[most]{tcolorbox}

% Boîte pour chaque rôle
\newtcolorbox{studentbox}[1][]{
  enhanced, breakable,
  colframe=studentOrange, colback=studentBg!30,
  coltitle=white, colbacktitle=studentOrange,
  fonttitle=\bfseries\large,
  title={#1},
  boxrule=1pt, arc=4pt,
  left=8pt, right=8pt, top=6pt, bottom=6pt,
  drop shadow={opacity=0.15},
}

\newtcolorbox{teacherbox}[1][]{
  enhanced, breakable,
  colframe=teacherGreen, colback=teacherBg!30,
  coltitle=white, colbacktitle=teacherGreen,
  fonttitle=\bfseries\large,
  title={#1},
  boxrule=1pt, arc=4pt,
  left=8pt, right=8pt, top=6pt, bottom=6pt,
  drop shadow={opacity=0.15},
}

\newtcolorbox{chefbox}[1][]{
  enhanced, breakable,
  colframe=chefBlue, colback=chefBg!30,
  coltitle=white, colbacktitle=chefBlue,
  fonttitle=\bfseries\large,
  title={#1},
  boxrule=1pt, arc=4pt,
  left=8pt, right=8pt, top=6pt, bottom=6pt,
  drop shadow={opacity=0.15},
}

% Boîte d'info / légende
\newtcolorbox{infobox}[1][]{
  enhanced,
  colframe=accentGold, colback=accentGold!8,
  coltitle=darkText, colbacktitle=accentGold!25,
  fonttitle=\bfseries,
  title={#1},
  boxrule=0.8pt, arc=3pt,
  left=6pt, right=6pt,
}

% ── En-têtes / pieds de page ──────────────────────────────
\usepackage{fancyhdr}
\setlength{\headheight}{14pt}
\pagestyle{fancy}
\fancyhf{}
\fancyhead[L]{\small\color{routeGray}\textit{User Stories — Projet PFE}}
\fancyhead[R]{\small\color{routeGray}\textit{\today}}
\fancyfoot[C]{\small\color{routeGray}\thepage}
\renewcommand{\headrulewidth}{0.4pt}
\renewcommand{\footrulewidth}{0pt}

% ── Hyperliens ─────────────────────────────────────────────
\usepackage[hidelinks]{hyperref}

% ── Commandes utilitaires ──────────────────────────────────
\newcommand{\route}[1]{\texorpdfstring{\texttt{\small\color{routeGray}#1}}{#1}}
\newcommand{\httpmethod}[1]{%
  \ifstrequal{#1}{GET}{\colorbox{blue!15}{\texttt{\bfseries\color{blue!70!black}GET}}}{}%
  \ifstrequal{#1}{POST}{\colorbox{green!18}{\texttt{\bfseries\color{green!60!black}POST}}}{}%
  \ifstrequal{#1}{PUT}{\colorbox{orange!18}{\texttt{\bfseries\color{orange!70!black}PUT}}}{}%
  \ifstrequal{#1}{DELETE}{\colorbox{red!15}{\texttt{\bfseries\color{red!70!black}DELETE}}}{}%
}
\newcommand{\usmark}{\color{accentGold}$\blacktriangleright$\;}
\newcommand{\cmark}{\textcolor{teacherGreen}{\ding{51}}}
\newcommand{\xmark}{\textcolor{red}{\ding{55}}}

% ── Titre ──────────────────────────────────────────────────
\usepackage{tikz}
\usetikzlibrary{shadows}

% ============================================================
\begin{document}
\hypersetup{pageanchor=false}

% ── PAGE DE TITRE ──────────────────────────────────────────
\begin{titlepage}
  \centering
  \vspace*{3cm}

  {\Huge\bfseries\color{headerBg} Cahier des User Stories}\\[0.6cm]
  {\Large\color{routeGray} Projet de Fin d'Études}\\[1.5cm]

  \begin{tikzpicture}
    \draw[headerBg, line width=2pt] (0,0) -- (12,0);
  \end{tikzpicture}\\[1.5cm]

  {\large\color{darkText}
    Analyse complète du \textbf{Backend} :\\[0.3cm]
    \textit{Rôles : Étudiant · Enseignant · Chef de Département}
  }\\[2cm]

  \begin{infobox}[\centering Modules analysés]
    \centering
    \small
    Authentification \textbullet\;
    Départements \textbullet\;
    Académique \textbullet\;
    Notes \textbullet\;
    Présence \textbullet\;
    Réclamations \textbullet\;
    Chat \textbullet\;
    Actualités \textbullet\;
    Événements \textbullet\;
    Stages \textbullet\;
    Utilisateurs
  \end{infobox}

  \vfill
  {\small\color{routeGray} Document généré automatiquement à partir de l'analyse du code source.}\\[0.5cm]
  {\small\color{routeGray} \today}
\end{titlepage}
\hypersetup{pageanchor=true}

% ── TABLE DES MATIÈRES ─────────────────────────────────────
\tableofcontents
\newpage

% ============================================================
\section{Introduction}
% ============================================================

Ce document présente les \textbf{user stories} (récits utilisateur) extraites de l'analyse complète du backend de l'application. Chaque user story décrit une fonctionnalité accessible à un rôle donné, accompagnée de la route API correspondante.

\subsection{Rôles du système}

Le système définit \textbf{cinq rôles} dans le modèle \texttt{User}, dont trois principaux :

\begin{center}
\begin{tabular}{llp{8cm}}
  \toprule
  \textbf{Rôle} & \textbf{Middleware} & \textbf{Description} \\
  \midrule
  \textbf{STUDENT}     & \texttt{protect}             & Accès limité à ses propres données personnelles \\
  \textbf{TEACHER}     & \texttt{protect + teacher}    & Gestion des notes et de la présence pour ses étudiants \\
  \textbf{CHEF\_DEPT}  & \texttt{protect + chefDept}   & Administration complète du département + capacités enseignant \\
  \textbf{ADMIN}       & \texttt{protect + admin}      & Super-utilisateur, passe tous les contrôles d'accès \\
  \textbf{PARTNER}     & \texttt{protect}              & Rôle défini mais sans routes spécifiques actuellement \\
  \bottomrule
\end{tabular}
\end{center}

\subsection{Légende des méthodes HTTP}

\begin{center}
  \httpmethod{GET} Lecture \qquad
  \httpmethod{POST} Création \qquad
  \httpmethod{PUT} Modification \qquad
  \httpmethod{DELETE} Suppression
\end{center}

\subsection{Hiérarchie des permissions}

\begin{infobox}[Héritage des droits]
  Le middleware \texttt{chefDept} autorise aussi le rôle \texttt{ADMIN}.
  Le middleware \texttt{teacher} autorise \texttt{TEACHER}, \texttt{CHEF\_DEPT} et \texttt{ADMIN}.
  Cela signifie :
  \[
    \textbf{STUDENT} \subset \textbf{TEACHER} \subset \textbf{CHEF\_DEPT} \subset \textbf{ADMIN}
  \]
  Chaque rôle supérieur hérite des capacités du rôle inférieur.
\end{infobox}

\newpage

% ============================================================
\section{User Stories — Étudiant (\texttt{STUDENT})}
% ============================================================

\begin{studentbox}[\Large\ding{56} Rôle : Étudiant]

% ─── Authentification ──────────────────────────────────────
\subsection*{Authentification}

\begin{itemize}[leftmargin=1.2cm, labelsep=0.4cm]
  \item[\usmark] \textbf{En tant qu'étudiant}, je peux \textbf{créer un compte} en fournissant mon nom, prénom, email et mot de passe.\\
    \httpmethod{POST} \route{/api/auth/register}

  \item[\usmark] \textbf{En tant qu'étudiant}, je peux \textbf{me connecter} avec mon email et mot de passe pour obtenir un jeton JWT.\\
    \httpmethod{POST} \route{/api/auth/login}

  \item[\usmark] \textbf{En tant qu'étudiant}, je peux \textbf{consulter mon profil} (nom, email, rôle, département, identifiant étudiant).\\
    \httpmethod{GET} \route{/api/auth/profile}
\end{itemize}

% ─── Emploi du temps ───────────────────────────────────────
\subsection*{Emploi du temps}

\begin{itemize}[leftmargin=1.2cm, labelsep=0.4cm]
  \item[\usmark] \textbf{En tant qu'étudiant}, je peux \textbf{consulter mon emploi du temps} personnalisé, résolu automatiquement à partir de la classe (\texttt{classId}) associée à mon profil.\\
    \httpmethod{GET} \route{/api/academic/schedule/student}
\end{itemize}

% ─── Notes ─────────────────────────────────────────────────
\subsection*{Notes}

\begin{itemize}[leftmargin=1.2cm, labelsep=0.4cm]
  \item[\usmark] \textbf{En tant qu'étudiant}, je peux \textbf{consulter mes notes}, filtrées par semestre. Le système restreint automatiquement la requête à mes propres notes uniquement.\\
    \httpmethod{GET} \route{/api/grades}
\end{itemize}

% ─── Présence ──────────────────────────────────────────────
\subsection*{Présence}

\begin{itemize}[leftmargin=1.2cm, labelsep=0.4cm]
  \item[\usmark] \textbf{En tant qu'étudiant}, je peux \textbf{consulter mes enregistrements de présence/absence}. Le système restreint la requête à mes propres enregistrements.\\
    \httpmethod{GET} \route{/api/attendance}
\end{itemize}

% ─── Réclamations ──────────────────────────────────────────
\subsection*{Réclamations de notes}

\begin{itemize}[leftmargin=1.2cm, labelsep=0.4cm]
  \item[\usmark] \textbf{En tant qu'étudiant}, je peux \textbf{déposer une réclamation} sur une de mes notes en précisant l'identifiant de la note (\texttt{gradeId}) et le motif de la réclamation.\\
    \textit{Contraintes : je ne peux réclamer que mes propres notes et je ne peux pas déposer deux réclamations sur la même note.}\\
    \httpmethod{POST} \route{/api/complaints}

  \item[\usmark] \textbf{En tant qu'étudiant}, je peux \textbf{consulter mes réclamations} et suivre leur statut (en attente, acceptée, rejetée).\\
    \httpmethod{GET} \route{/api/complaints}
\end{itemize}

% ─── Chat ──────────────────────────────────────────────────
\subsection*{Messagerie instantanée (Chat)}

\begin{itemize}[leftmargin=1.2cm, labelsep=0.4cm]
  \item[\usmark] \textbf{En tant qu'étudiant}, je peux \textbf{accéder au salon de chat général} et au \textbf{salon de ma classe}.\\
    \httpmethod{GET} \route{/api/chat/rooms}

  \item[\usmark] \textbf{En tant qu'étudiant}, je peux \textbf{consulter l'historique} des messages d'un salon.\\
    \httpmethod{GET} \route{/api/chat/history/:room}

  \item[\usmark] \textbf{En tant qu'étudiant}, je peux \textbf{envoyer des messages} en temps réel, rejoindre/quitter des salons, et voir les indicateurs de frappe via \textbf{WebSocket}.
\end{itemize}

% ─── Contenu public ────────────────────────────────────────
\subsection*{Contenu public (sans authentification requise)}

\begin{itemize}[leftmargin=1.2cm, labelsep=0.4cm]
  \item[\usmark] \textbf{En tant qu'étudiant}, je peux \textbf{consulter la liste des actualités} et lire un article en détail.\\
    \httpmethod{GET} \route{/api/news} \quad \httpmethod{GET} \route{/api/news/:id}

  \item[\usmark] \textbf{En tant qu'étudiant}, je peux \textbf{consulter la liste des événements} et voir les détails d'un événement.\\
    \httpmethod{GET} \route{/api/events} \quad \httpmethod{GET} \route{/api/events/:id}

  \item[\usmark] \textbf{En tant qu'étudiant}, je peux \textbf{consulter les offres de stages} ouvertes et voir les détails d'une offre.\\
    \httpmethod{GET} \route{/api/stages} \quad \httpmethod{GET} \route{/api/stages/:id}
\end{itemize}

\end{studentbox}

\newpage

% ============================================================
\section{User Stories — Enseignant (\texttt{TEACHER})}
% ============================================================

\begin{teacherbox}[\Large\ding{56} Rôle : Enseignant]

\begin{infobox}[Héritage]
  L'enseignant \textbf{hérite de toutes les capacités de lecture} de l'étudiant (emploi du temps, contenu public, chat) et acquiert des \textbf{capacités de gestion} supplémentaires.
\end{infobox}

% ─── Emploi du temps ───────────────────────────────────────
\subsection*{Emploi du temps}

\begin{itemize}[leftmargin=1.2cm, labelsep=0.4cm]
  \item[\usmark] \textbf{En tant qu'enseignant}, je peux \textbf{consulter mon emploi du temps personnel}, résolu par correspondance de mon nom et prénom dans les séances.\\
    \httpmethod{GET} \route{/api/academic/schedule/teacher}
\end{itemize}

% ─── Cours & Classes ──────────────────────────────────────
\subsection*{Mes cours et classes}

\begin{itemize}[leftmargin=1.2cm, labelsep=0.4cm]
  \item[\usmark] \textbf{En tant qu'enseignant}, je peux \textbf{consulter la liste des cours que j'enseigne}, déduits automatiquement des séances programmées.\\
    \httpmethod{GET} \route{/api/academic/teacher/courses}

  \item[\usmark] \textbf{En tant qu'enseignant}, je peux \textbf{consulter la liste des classes auxquelles j'enseigne}, avec les matières associées et le nombre de séances.\\
    \httpmethod{GET} \route{/api/academic/teacher/classes}
\end{itemize}

% ─── Notes ─────────────────────────────────────────────────
\subsection*{Gestion des notes}

\begin{itemize}[leftmargin=1.2cm, labelsep=0.4cm]
  \item[\usmark] \textbf{En tant qu'enseignant}, je peux \textbf{consulter les notes} que j'ai attribuées, filtrées par \texttt{studentId} et/ou \texttt{semester}.\\
    \httpmethod{GET} \route{/api/grades}

  \item[\usmark] \textbf{En tant qu'enseignant}, je peux \textbf{ajouter une note} pour un étudiant (matière, score, coefficient, type, semestre).\\
    \httpmethod{POST} \route{/api/grades}

  \item[\usmark] \textbf{En tant qu'enseignant}, je peux \textbf{supprimer une note} que j'ai créée.\\
    \httpmethod{DELETE} \route{/api/grades/:id}

  \item[\usmark] \textbf{En tant qu'enseignant}, je peux \textbf{importer des notes en masse} depuis un fichier Excel (colonnes : \texttt{studentId}, \texttt{score}) pour toute une classe.\\
    \httpmethod{POST} \route{/api/grades/bulk-upload}
\end{itemize}

% ─── Présence ──────────────────────────────────────────────
\subsection*{Gestion de la présence}

\begin{itemize}[leftmargin=1.2cm, labelsep=0.4cm]
  \item[\usmark] \textbf{En tant qu'enseignant}, je peux \textbf{consulter les enregistrements de présence} de mes séances.\\
    \httpmethod{GET} \route{/api/attendance}

  \item[\usmark] \textbf{En tant qu'enseignant}, je peux \textbf{marquer la présence d'un étudiant} individuellement (présent, absent, en retard).\\
    \httpmethod{POST} \route{/api/attendance}

  \item[\usmark] \textbf{En tant qu'enseignant}, je peux \textbf{marquer la présence en masse} pour toute une séance de classe.\\
    \httpmethod{POST} \route{/api/attendance/bulk}

  \item[\usmark] \textbf{En tant qu'enseignant}, je peux \textbf{modifier un enregistrement de présence} (statut, justification).\\
    \httpmethod{PUT} \route{/api/attendance/:id}

  \item[\usmark] \textbf{En tant qu'enseignant}, je peux \textbf{supprimer un enregistrement de présence}.\\
    \httpmethod{DELETE} \route{/api/attendance/:id}

  \item[\usmark] \textbf{En tant qu'enseignant}, je peux \textbf{récupérer la liste des étudiants d'une classe} par nom de classe, pour faciliter la prise de présence.\\
    \httpmethod{GET} \route{/api/attendance/students-by-class}
\end{itemize}

% ─── Utilisateurs ──────────────────────────────────────────
\subsection*{Liste des utilisateurs}

\begin{itemize}[leftmargin=1.2cm, labelsep=0.4cm]
  \item[\usmark] \textbf{En tant qu'enseignant}, je peux \textbf{consulter la liste des étudiants}. Le système restreint automatiquement les résultats au rôle \texttt{STUDENT}.\\
    \httpmethod{GET} \route{/api/users}
\end{itemize}

% ─── Chat ──────────────────────────────────────────────────
\subsection*{Messagerie instantanée}

\begin{itemize}[leftmargin=1.2cm, labelsep=0.4cm]
  \item[\usmark] \textbf{En tant qu'enseignant}, je peux \textbf{accéder au salon général} et aux \textbf{salons de toutes les classes auxquelles j'enseigne}.\\
    \httpmethod{GET} \route{/api/chat/rooms}

  \item[\usmark] \textbf{En tant qu'enseignant}, je peux envoyer des messages et interagir en temps réel via WebSocket.
\end{itemize}

\end{teacherbox}

\newpage

% ============================================================
\section{User Stories — Chef de Département (\texttt{CHEF\_DEPT})}
% ============================================================

\begin{chefbox}[\Large\ding{56} Rôle : Chef de Département]

\begin{infobox}[Héritage]
  Le Chef de Département \textbf{hérite de toutes les capacités de l'Enseignant} (notes, présence, cours, etc.) et acquiert l'\textbf{administration complète du département}.
\end{infobox}

% ─── Département ───────────────────────────────────────────
\subsection*{Gestion du département}

\begin{itemize}[leftmargin=1.2cm, labelsep=0.4cm]
  \item[\usmark] \textbf{En tant que Chef}, je peux \textbf{consulter tous les départements} avec le nombre d'enseignants.\\
    \httpmethod{GET} \route{/api/departments}

  \item[\usmark] \textbf{En tant que Chef}, je peux \textbf{consulter un département en détail} (tous les enseignants, classes, etc.).\\
    \httpmethod{GET} \route{/api/departments/:id}

  \item[\usmark] \textbf{En tant que Chef}, je peux \textbf{consulter mon propre département}, identifié par mon adresse email.\\
    \httpmethod{GET} \route{/api/departments/my}
\end{itemize}

% ─── Enseignants ───────────────────────────────────────────
\subsection*{Gestion des enseignants}

\begin{itemize}[leftmargin=1.2cm, labelsep=0.4cm]
  \item[\usmark] \textbf{En tant que Chef}, je peux \textbf{ajouter un enseignant} à mon département (nom, prénom, email, grade, spécialisation).\\
    \httpmethod{POST} \route{/api/departments/my/teachers}

  \item[\usmark] \textbf{En tant que Chef}, je peux \textbf{modifier les informations} d'un enseignant de mon département.\\
    \httpmethod{PUT} \route{/api/departments/my/teachers/:teacherId}

  \item[\usmark] \textbf{En tant que Chef}, je peux \textbf{supprimer un enseignant} de mon département.\\
    \httpmethod{DELETE} \route{/api/departments/my/teachers/:teacherId}
\end{itemize}

% ─── Classes ───────────────────────────────────────────────
\subsection*{Gestion des classes}

\begin{itemize}[leftmargin=1.2cm, labelsep=0.4cm]
  \item[\usmark] \textbf{En tant que Chef}, je peux \textbf{créer une nouvelle classe} (nom, niveau, parcours, année académique).\\
    \httpmethod{POST} \route{/api/departments/my/classes}

  \item[\usmark] \textbf{En tant que Chef}, je peux \textbf{modifier une classe} existante.\\
    \httpmethod{PUT} \route{/api/departments/my/classes/:classId}

  \item[\usmark] \textbf{En tant que Chef}, je peux \textbf{supprimer une classe}.\\
    \httpmethod{DELETE} \route{/api/departments/my/classes/:classId}

  \item[\usmark] \textbf{En tant que Chef}, je peux \textbf{importer des étudiants en masse} dans une classe à partir d'une liste (nom, prénom, email, numéro d'inscription). Un mot de passe par défaut (\texttt{iset123}) est attribué automatiquement.\\
    \httpmethod{POST} \route{/api/departments/my/classes/:classId/students/bulk}
\end{itemize}

% ─── Étudiants ─────────────────────────────────────────────
\subsection*{Gestion des étudiants}

\begin{itemize}[leftmargin=1.2cm, labelsep=0.4cm]
  \item[\usmark] \textbf{En tant que Chef}, je peux \textbf{consulter tous les étudiants} de mon département.\\
    \httpmethod{GET} \route{/api/academic/students}

  \item[\usmark] \textbf{En tant que Chef}, je peux \textbf{ajouter un étudiant} à mon département avec mot de passe par défaut.\\
    \httpmethod{POST} \route{/api/academic/students}

  \item[\usmark] \textbf{En tant que Chef}, je peux \textbf{modifier le profil d'un étudiant} (nom, email, classe, numéro d'inscription).\\
    \httpmethod{PUT} \route{/api/academic/students/:id}

  \item[\usmark] \textbf{En tant que Chef}, je peux \textbf{supprimer un étudiant} de mon département.\\
    \httpmethod{DELETE} \route{/api/academic/students/:id}
\end{itemize}

% ─── Cours ─────────────────────────────────────────────────
\subsection*{Gestion des cours}

\begin{itemize}[leftmargin=1.2cm, labelsep=0.4cm]
  \item[\usmark] \textbf{En tant que Chef}, je peux \textbf{consulter tous les cours} de mon département.\\
    \httpmethod{GET} \route{/api/academic/courses}

  \item[\usmark] \textbf{En tant que Chef}, je peux \textbf{ajouter un cours} (nom, code, semestre, niveau, volume horaire).\\
    \httpmethod{POST} \route{/api/academic/courses}

  \item[\usmark] \textbf{En tant que Chef}, je peux \textbf{modifier un cours} existant.\\
    \httpmethod{PUT} \route{/api/academic/courses/:id}

  \item[\usmark] \textbf{En tant que Chef}, je peux \textbf{supprimer un cours}.\\
    \httpmethod{DELETE} \route{/api/academic/courses/:id}
\end{itemize}

% ─── Emploi du temps ───────────────────────────────────────
\subsection*{Gestion de l'emploi du temps}

\begin{itemize}[leftmargin=1.2cm, labelsep=0.4cm]
  \item[\usmark] \textbf{En tant que Chef}, je peux \textbf{consulter l'emploi du temps} de n'importe quelle classe de mon département.\\
    \httpmethod{GET} \route{/api/academic/schedule}

  \item[\usmark] \textbf{En tant que Chef}, je peux \textbf{consulter toutes les salles} disponibles.\\
    \httpmethod{GET} \route{/api/academic/rooms}

  \item[\usmark] \textbf{En tant que Chef}, je peux \textbf{créer une séance} (affecter un cours, enseignant, salle, créneau horaire à une classe).\\
    \httpmethod{POST} \route{/api/academic/sessions}

  \item[\usmark] \textbf{En tant que Chef}, je peux \textbf{modifier une séance} programmée.\\
    \httpmethod{PUT} \route{/api/academic/sessions/:id}

  \item[\usmark] \textbf{En tant que Chef}, je peux \textbf{supprimer une séance}.\\
    \httpmethod{DELETE} \route{/api/academic/sessions/:id}
\end{itemize}

% ─── Réclamations ──────────────────────────────────────────
\subsection*{Gestion des réclamations}

\begin{itemize}[leftmargin=1.2cm, labelsep=0.4cm]
  \item[\usmark] \textbf{En tant que Chef}, je peux \textbf{consulter toutes les réclamations} de tous les étudiants (pas seulement les miens).\\
    \httpmethod{GET} \route{/api/complaints}

  \item[\usmark] \textbf{En tant que Chef}, je peux \textbf{traiter une réclamation} en l'\textbf{acceptant} ou la \textbf{rejetant}, avec un message de réponse optionnel.\\
    \httpmethod{PUT} \route{/api/complaints/:id/resolve}
\end{itemize}

% ─── Actualités ────────────────────────────────────────────
\subsection*{Gestion des actualités}

\begin{itemize}[leftmargin=1.2cm, labelsep=0.4cm]
  \item[\usmark] \textbf{En tant que Chef}, je peux \textbf{créer un article d'actualité} avec image et document en pièce jointe.\\
    \httpmethod{POST} \route{/api/news}

  \item[\usmark] \textbf{En tant que Chef}, je peux \textbf{modifier un article} existant.\\
    \httpmethod{PUT} \route{/api/news/:id}

  \item[\usmark] \textbf{En tant que Chef}, je peux \textbf{supprimer un article}.\\
    \httpmethod{DELETE} \route{/api/news/:id}
\end{itemize}

% ─── Événements ────────────────────────────────────────────
\subsection*{Gestion des événements}

\begin{itemize}[leftmargin=1.2cm, labelsep=0.4cm]
  \item[\usmark] \textbf{En tant que Chef}, je peux \textbf{créer un événement} (titre, description, dates, lieu, audience cible) avec pièces jointes.\\
    \httpmethod{POST} \route{/api/events}

  \item[\usmark] \textbf{En tant que Chef}, je peux \textbf{modifier un événement}.\\
    \httpmethod{PUT} \route{/api/events/:id}

  \item[\usmark] \textbf{En tant que Chef}, je peux \textbf{supprimer un événement}.\\
    \httpmethod{DELETE} \route{/api/events/:id}
\end{itemize}

% ─── Stages ────────────────────────────────────────────────
\subsection*{Gestion des offres de stages}

\begin{itemize}[leftmargin=1.2cm, labelsep=0.4cm]
  \item[\usmark] \textbf{En tant que Chef}, je peux \textbf{créer une offre de stage} (titre, entreprise, lieu, durée, type, description, email de contact, date limite).\\
    \httpmethod{POST} \route{/api/stages}

  \item[\usmark] \textbf{En tant que Chef}, je peux \textbf{modifier une offre de stage}.\\
    \httpmethod{PUT} \route{/api/stages/:id}

  \item[\usmark] \textbf{En tant que Chef}, je peux \textbf{supprimer une offre de stage}.\\
    \httpmethod{DELETE} \route{/api/stages/:id}
\end{itemize}

% ─── Chat ──────────────────────────────────────────────────
\subsection*{Messagerie instantanée}

\begin{itemize}[leftmargin=1.2cm, labelsep=0.4cm]
  \item[\usmark] \textbf{En tant que Chef}, je peux \textbf{accéder au salon général} et aux \textbf{salons de toutes les classes de mon département}.\\
    \httpmethod{GET} \route{/api/chat/rooms}
\end{itemize}

\end{chefbox}

\newpage

% ============================================================
\section{Matrice récapitulative des fonctionnalités}
% ============================================================

Le tableau ci-dessous résume les accès de chaque rôle aux fonctionnalités principales du système.

\vspace{0.5cm}

\begin{center}
\small
\begin{tabular}{>{\raggedright\arraybackslash}p{5.5cm}ccc}
  \toprule
  \textbf{Fonctionnalité} & \textbf{Étudiant} & \textbf{Enseignant} & \textbf{Chef Dept} \\
  \midrule
  Inscription / Connexion / Profil & \cmark & \cmark & \cmark \\
  Consulter son emploi du temps & \cmark & \cmark & \cmark \\
  Consulter ses notes & \cmark\; \footnotesize{(lecture)} & \cmark\; \footnotesize{+ écriture} & \cmark\; \footnotesize{(hérité)} \\
  Ajouter / Supprimer des notes & \xmark & \cmark & \cmark\; \footnotesize{(hérité)} \\
  Import en masse de notes (Excel) & \xmark & \cmark & \cmark\; \footnotesize{(hérité)} \\
  Consulter la présence & \cmark\; \footnotesize{(lecture)} & \cmark\; \footnotesize{+ écriture} & \cmark\; \footnotesize{(hérité)} \\
  Marquer / Modifier / Supprimer la présence & \xmark & \cmark & \cmark\; \footnotesize{(hérité)} \\
  Présence en masse & \xmark & \cmark & \cmark\; \footnotesize{(hérité)} \\
  Déposer une réclamation de note & \cmark & \xmark & \xmark \\
  Traiter les réclamations & \xmark & \xmark & \cmark \\
  Gérer les enseignants & \xmark & \xmark & \cmark\; \footnotesize{CRUD} \\
  Gérer les classes & \xmark & \xmark & \cmark\; \footnotesize{CRUD} \\
  Gérer les étudiants & \xmark & \xmark & \cmark\; \footnotesize{CRUD + mass} \\
  Gérer les cours (modules) & \xmark & \xmark & \cmark\; \footnotesize{CRUD} \\
  Gérer les séances (emploi du temps) & \xmark & \xmark & \cmark\; \footnotesize{CRUD} \\
  Gérer les actualités & \xmark & \xmark & \cmark\; \footnotesize{CRUD} \\
  Gérer les événements & \xmark & \xmark & \cmark\; \footnotesize{CRUD} \\
  Gérer les offres de stages & \xmark & \xmark & \cmark\; \footnotesize{CRUD} \\
  Chat (salon général) & \cmark & \cmark & \cmark \\
  Chat (salons de classe) & \cmark\; \footnotesize{(sa classe)} & \cmark\; \footnotesize{(ses classes)} & \cmark\; \footnotesize{(tout le dept)} \\
  Parcourir actualités / événements / stages & \cmark\; \footnotesize{Public} & \cmark\; \footnotesize{Public} & \cmark\; \footnotesize{Public} \\
  Consulter la liste des utilisateurs & \xmark & \cmark\; \footnotesize{(étudiants)} & \cmark\; \footnotesize{(tous)} \\
  \bottomrule
\end{tabular}
\end{center}

\newpage

% ============================================================
\section{Annexe — Routes API complètes}
% ============================================================

\subsection{Authentification — \route{/api/auth}}

\begin{center}
\begin{tabular}{|l|l|p{6cm}|l|}
  \hline
  \rowcolor{headerBg}
  \textcolor{white}{\textbf{Méthode}} & \textcolor{white}{\textbf{Route}} & \textcolor{white}{\textbf{Description}} & \textcolor{white}{\textbf{Accès}} \\
  \hline
  \httpmethod{POST} & \texttt{/register} & Créer un compte & Public \\
  \httpmethod{POST} & \texttt{/login} & Se connecter & Public \\
  \httpmethod{GET}  & \texttt{/profile} & Voir son profil & Authentifié \\
  \hline
\end{tabular}
\end{center}

\subsection{Académique — \route{/api/academic}}

\begin{center}
\small
\begin{tabular}{|l|l|p{6cm}|l|}
  \hline
  \rowcolor{headerBg}
  \textcolor{white}{\textbf{Méthode}} & \textcolor{white}{\textbf{Route}} & \textcolor{white}{\textbf{Description}} & \textcolor{white}{\textbf{Accès}} \\
  \hline
  \httpmethod{GET} & \texttt{/schedule/student} & Emploi du temps étudiant & Étudiant \\
  \httpmethod{GET} & \texttt{/schedule/teacher} & Emploi du temps enseignant & Enseignant \\
  \httpmethod{GET} & \texttt{/teacher/courses} & Cours de l'enseignant & Enseignant \\
  \httpmethod{GET} & \texttt{/teacher/classes} & Classes de l'enseignant & Enseignant \\
  \httpmethod{GET} & \texttt{/courses} & Cours du département & Chef \\
  \httpmethod{GET} & \texttt{/rooms} & Salles disponibles & Chef \\
  \httpmethod{GET} & \texttt{/schedule} & Emploi du temps (par classe) & Chef \\
  \httpmethod{GET} & \texttt{/students} & Étudiants du département & Chef \\
  \httpmethod{POST} & \texttt{/students} & Ajouter un étudiant & Chef \\
  \httpmethod{PUT} & \texttt{/students/:id} & Modifier un étudiant & Chef \\
  \httpmethod{DELETE} & \texttt{/students/:id} & Supprimer un étudiant & Chef \\
  \httpmethod{POST} & \texttt{/courses} & Ajouter un cours & Chef \\
  \httpmethod{PUT} & \texttt{/courses/:id} & Modifier un cours & Chef \\
  \httpmethod{DELETE} & \texttt{/courses/:id} & Supprimer un cours & Chef \\
  \httpmethod{POST} & \texttt{/sessions} & Créer une séance & Chef \\
  \httpmethod{PUT} & \texttt{/sessions/:id} & Modifier une séance & Chef \\
  \httpmethod{DELETE} & \texttt{/sessions/:id} & Supprimer une séance & Chef \\
  \hline
\end{tabular}
\end{center}

\subsection{Départements — \route{/api/departments}}

\begin{center}
\small
\begin{tabular}{l>{\raggedright\arraybackslash}p{5.6cm}>{\raggedright\arraybackslash}p{4.4cm}l}
  \toprule
  \textbf{Méthode} & \textbf{Route} & \textbf{Description} & \textbf{Accès} \\
  \midrule
  \httpmethod{GET} & \texttt{/} & Liste des départements & Authentifié \\
  \httpmethod{GET} & \texttt{/:id} & Détails d'un département & Authentifié \\
  \httpmethod{GET} & \texttt{/my} & Mon département & Authentifié \\
  \httpmethod{POST} & \texttt{/my/teachers} & Ajouter un enseignant & Chef \\
  \httpmethod{PUT} & \texttt{/my/teachers/:id} & Modifier un enseignant & Chef \\
  \httpmethod{DELETE} & \texttt{/my/teachers/:id} & Supprimer un enseignant & Chef \\
  \httpmethod{POST} & \texttt{/my/classes} & Créer une classe & Chef \\
  \httpmethod{PUT} & \texttt{/my/classes/:id} & Modifier une classe & Chef \\
  \httpmethod{DELETE} & \texttt{/my/classes/:id} & Supprimer une classe & Chef \\
  \httpmethod{POST} & \texttt{/my/classes/:id/students/}\newline\texttt{bulk} & Import en masse & Chef \\
  \bottomrule
\end{tabular}
\end{center}

\subsection{Notes — \route{/api/grades}}

\begin{center}
\begin{tabular}{l>{\raggedright\arraybackslash}p{5.6cm}>{\raggedright\arraybackslash}p{4.4cm}l}
  \toprule
  \textbf{Méthode} & \textbf{Route} & \textbf{Description} & \textbf{Accès} \\
  \midrule
  \httpmethod{GET} & \texttt{/} & Consulter les notes & Authentifié \\
  \httpmethod{POST} & \texttt{/} & Ajouter une note & Enseignant \\
  \httpmethod{POST} & \texttt{/bulk-upload} & Import Excel & Enseignant \\
  \httpmethod{DELETE} & \texttt{/:id} & Supprimer une note & Enseignant \\
  \bottomrule
\end{tabular}
\end{center}

\subsection{Présence — \route{/api/attendance}}

\begin{center}
\begin{tabular}{l>{\raggedright\arraybackslash}p{5.6cm}>{\raggedright\arraybackslash}p{4.4cm}l}
  \toprule
  \textbf{Méthode} & \textbf{Route} & \textbf{Description} & \textbf{Accès} \\
  \midrule
  \httpmethod{GET} & \texttt{/} & Consulter la présence & Authentifié \\
  \httpmethod{POST} & \texttt{/} & Marquer (individuel) & Enseignant \\
  \httpmethod{POST} & \texttt{/bulk} & Marquer (en masse) & Enseignant \\
  \httpmethod{GET} & \texttt{/students-by-class} & Étudiants d'une classe & Enseignant \\
  \httpmethod{PUT} & \texttt{/:id} & Modifier un enregistrement & Enseignant \\
  \httpmethod{DELETE} & \texttt{/:id} & Supprimer un enregistrement & Enseignant \\
  \bottomrule
\end{tabular}
\end{center}

\subsection{Réclamations — \route{/api/complaints}}

\begin{center}
\begin{tabular}{l>{\raggedright\arraybackslash}p{5.6cm}>{\raggedright\arraybackslash}p{4.4cm}l}
  \toprule
  \textbf{Méthode} & \textbf{Route} & \textbf{Description} & \textbf{Accès} \\
  \midrule
  \httpmethod{GET} & \texttt{/} & Consulter les réclamations & Authentifié \\
  \httpmethod{POST} & \texttt{/} & Déposer une réclamation & Authentifié \\
  \httpmethod{PUT} & \texttt{/:id/resolve} & Traiter une réclamation & Chef \\
  \bottomrule
\end{tabular}
\end{center}

\subsection{Chat — \route{/api/chat}}

\begin{center}
\begin{tabular}{l>{\raggedright\arraybackslash}p{5.6cm}>{\raggedright\arraybackslash}p{4.4cm}l}
  \toprule
  \textbf{Méthode} & \textbf{Route} & \textbf{Description} & \textbf{Accès} \\
  \midrule
  \httpmethod{GET} & \texttt{/rooms} & Salons disponibles & Authentifié \\
  \httpmethod{GET} & \texttt{/history/:room} & Historique d'un salon & Authentifié \\
  \bottomrule
\end{tabular}
\end{center}

\subsection{Actualités — \route{/api/news}}

\begin{center}
\begin{tabular}{l>{\raggedright\arraybackslash}p{5.6cm}>{\raggedright\arraybackslash}p{4.4cm}l}
  \toprule
  \textbf{Méthode} & \textbf{Route} & \textbf{Description} & \textbf{Accès} \\
  \midrule
  \httpmethod{GET} & \texttt{/} & Liste des actualités & Public \\
  \httpmethod{GET} & \texttt{/:id} & Détails d'un article & Public \\
  \httpmethod{POST} & \texttt{/} & Créer un article & Chef \\
  \httpmethod{PUT} & \texttt{/:id} & Modifier un article & Chef \\
  \httpmethod{DELETE} & \texttt{/:id} & Supprimer un article & Chef \\
  \bottomrule
\end{tabular}
\end{center}

\subsection{Événements — \route{/api/events}}

\begin{center}
\begin{tabular}{l>{\raggedright\arraybackslash}p{5.6cm}>{\raggedright\arraybackslash}p{4.4cm}l}
  \toprule
  \textbf{Méthode} & \textbf{Route} & \textbf{Description} & \textbf{Accès} \\
  \midrule
  \httpmethod{GET} & \texttt{/} & Liste des événements & Public \\
  \httpmethod{GET} & \texttt{/:id} & Détails d'un événement & Public \\
  \httpmethod{POST} & \texttt{/} & Créer un événement & Chef \\
  \httpmethod{PUT} & \texttt{/:id} & Modifier un événement & Chef \\
  \httpmethod{DELETE} & \texttt{/:id} & Supprimer un événement & Chef \\
  \bottomrule
\end{tabular}
\end{center}

\subsection{Stages — \route{/api/stages}}

\begin{center}
\begin{tabular}{l>{\raggedright\arraybackslash}p{5.6cm}>{\raggedright\arraybackslash}p{4.4cm}l}
  \toprule
  \textbf{Méthode} & \textbf{Route} & \textbf{Description} & \textbf{Accès} \\
  \midrule
  \httpmethod{GET} & \texttt{/} & Liste des offres & Public \\
  \httpmethod{GET} & \texttt{/:id} & Détails d'une offre & Public \\
  \httpmethod{POST} & \texttt{/} & Créer une offre & Chef \\
  \httpmethod{PUT} & \texttt{/:id} & Modifier une offre & Chef \\
  \httpmethod{DELETE} & \texttt{/:id} & Supprimer une offre & Chef \\
  \bottomrule
\end{tabular}
\end{center}

\subsection{Utilisateurs — \route{/api/users}}

\begin{center}
\begin{tabular}{l>{\raggedright\arraybackslash}p{5.6cm}>{\raggedright\arraybackslash}p{4.4cm}l}
  \toprule
  \textbf{Méthode} & \textbf{Route} & \textbf{Description} & \textbf{Accès} \\
  \midrule
  \httpmethod{GET} & \texttt{/} & Liste des utilisateurs & Authentifié \\
  \bottomrule
\end{tabular}
\end{center}

\end{document}
